\subsection{Analytical Integrated Assessment Models}

\textcite{NordhausWilliam1979} sparked an entire field of research into IAMs, including his own Dynamic Integrated Climate-Economy (DICE) model \parencite{NordhausWilliamD2011, NordhausWilliam2018PaUa, NordhausWilliam2023}. Simply put, there are two main branches of IAMs: numerical IAMs, often used in policy advising, and analytical IAMs, which offer analytical insights into important quantities such as the social cost of carbon (SCC). The IAM used in this thesis is the Analytical Climate Economy (ACE) model, presented in \textcite{TraegerChristian2023ACE}. Therefore, the literature review will concentrate on those aspects of the literature that pertain directly to ACE\footnote{For reviews on numerical IAMs, see \textcite{BosettiValentina2021, MetcalfStock2017}.}. 

\textcite{Heal1984} marks one of the first efforts to develop analytical frameworks for IAMs and focused on policy design under uncertainty, but was nonquantitative. Later came contributions such as \textcite{SchauerMichael1995}, who developed an explicit model of the SCC under uncertainty and \textcite{Pizer1999}, who did a log-linearization of a DICE-based IAM. Other contributions were the linear quadratic models in \textcite{HoelKarp2002, KarpZhang2012}, but they did not model production- or energy sectors.

\textcite{Golosov2014etal} made significant progress in the field by specifying a Dynamic Stochastic General Equilibrium (DSGE) model with an energy sector and modeled the production response to emissions. This was an extension of non-renewable resource models like the one found in \textcite{DasguptaHeal1974}. One limitation of \textcite{Golosov2014etal} was that they assumed the production impulse response to emissions to be immediate and, therefore, also implicitly assumed an immediate response in temperatures. However, this is inconsistent with the observation that the oceans contribute to warming over centuries \parencite{Abraham2022ocean_resp}. \textcite{GerlaghLiski2018} extended the climate system representation to better model the persistence and delay from climate impacts.

Subsequent developments in Analytical Integrated Assessment Models (AIAMs), inspired by the work of \textcite{Golosov2014etal}, leveraged the property that the models that can be transformed to a system linear in the states can be solved analytically, as pointed out in \textcite{Karp2017}. A few notable applications of this approach include an application to international cooperation and intergenerational altruism in games that affect carbon levels \parencite{Karp2017} and a non-cooperative multi-regional setting as in \textcite{Hambel2021non-coop-world}.

\textcite{TraegerChristian2023ACE} presented the Analytical Climate Economy (ACE) model, which addressed four key limitations in the literature. Firstly, ACE incorporates more realistic temperature dynamics by explicitly modeling the non-linear relationship between CO2 concentrations and temperature changes. Second, he improved upon the overly simplistic full-depreciation assumptions prevalent in \textcite{Golosov2014etal}. Third, by including delayed responses between emissions and their corresponding economic and climate impacts. Finally, ACE allowed for variable degrees of risk aversion distinct from intertemporal substitution elasticity, enhancing its utility in analyzing policy implications under uncertainty.

The following thesis applies a state-of-the-art deep learning technique called \enquote{Deep Equilibrium Nets}, which is presented in the paper by \textcite{Azinovic2022DEQN}, to the analytical structure of ACE \parencite{TraegerChristian2023ACE}. The research work makes two important contributions to this field. Firstly, it represents one of the earliest applications to IAMs. Before this work, only two studies had applied this method to models based on DICE, as seen in \textcite{Folini2024ClimateInCE, friedl2023_deepUQ}. This positions the present thesis at the front of deep learning applications to IAMs.

Secondly, by mapping DEQNs onto ACE, which is general in its structure, the thesis not only tests the quality and adaptability of the method but should also facilitate broader applications. Furthermore, DEQNs were introduced to overcome the curse of dimensionality \parencite{Bellman1961}, which is the exact problem IAMs face when incorporating more rich and complex features. Given that this application to ACE is generalizable, it paves the way for developing more complex IAMs.